\section{Paso 5: ejecuciones piloto}
\label{sec:paso5-pilotos}
% Origen: 02-modelo-anylogic.md §7 (primer piloto, instrumentación, elección de cola) y SIM.md (log del 24/09 y del 05/10).

Los pilotos tenían tres objetivos: comprobar en el motor de AnyLogic lo que el verificador comprueba en
el arnés de prueba, medir el tiempo de cómputo y encontrar fallas antes de producir. Encontraron tres
fallas que habrían invalidado los resultados. Esta sección las presenta en el orden en que aparecieron.

\subsection{Demostración corta}

El primer piloto usó una demostración sintética: seis oleadas de 80 pasajeros separadas por 30~s, con 3~s
de servicio. Esos valores se eligieron para forzar una cola visible y no representan a Constitución. Con
el corte inicial a los 600~s quedaban peatones dentro del vestíbulo. Por eso se separó el corte de las
métricas (600~s) del fin técnico de la corrida (900~s). Con ese cambio, \EO{} y \EI{} drenaron los 480
pasajeros y terminaron con la cola vacía.

\subsection{Primer piloto de la franja completa: el sistema no drenaba}

El 24/09/2026 se corrió la franja de 07:00 a 09:30 con valores de prueba:
\begin{itemize}
  \item 80~\% de la demanda en tren;
  \item 60~s de demora;
  \item 120~s de descarga;
  \item servicio triangular (2; 3; 5)~s.
\end{itemize}
La corrida de \EO{} generó \num{17430} pasajeros, pero no drenó. A las 09:30 habían salido \num{12600}, y
al llegar al tope de 5~horas de modelo de la edición PLE, \num{14232}. La utilización de los molinetes
era de apenas 23~\%, mientras la cola llegaba a 497 personas. El cuello de botella no era la validación,
sino el movimiento de los peatones. La corrida tardó unos 43 minutos. La de \EI{} se cortó a los 787~s
de modelo sin un error visible.

\textbf{Instrumentación.} Las métricas medían desde la entrada a la cola, por lo que no veían a quien
todavía no había llegado a ella. Se agregaron contadores de peatones en el hall (ingresados que no
llegaron a un molinete) y en la zona paga, y una línea de progreso cada 15 minutos simulados. Así se vio
que, al tope de las 5~horas, unos \num{2700} peatones estaban trabados en el hall y solo 497 en las colas.
La memoria de cada experimento se subió de 512--\num{1024}~MB a \num{4096}~MB.

\textbf{Una hipótesis equivocada.} El atasco se atribuyó primero a la orientación de las colas, y se
invirtieron. La comparación posterior con el ejemplo oficial \emph{Subway Entrance Hall} mostró que la
orientación original era la correcta: en AnyLogic, la cabeza de la cola es el último punto de la línea.
La inversión se deshizo. El piloto que no drenaba tenía las colas bien orientadas, así que la causa
seguía sin encontrarse.

\subsection{La causa: la elección de cola no veía la fila}

Un contador agregado a la función de elección mostró que el 100~\% de las elecciones se hacían sin tiempo
de servicio asignado: \num{18545} de \num{18545}. AnyLogic evalúa la elección de cola de
\texttt{PedService} \emph{antes} de la acción de salida de \texttt{PedSource}, que es donde se sorteaba el
tiempo de validación de cada pasajero. El costo de cada cola se calculaba como pasajeros en la fila por
servicio asignado, y valía cero. La fila no pesaba y cada pasajero iba al molinete más alineado con su
entrada. Los molinetes 2 a 8, 19 y 20 quedaban sin uso, las colas de los demás desbordaban hacia el acceso
y 432 personas quedaban trabadas en el hall.

La corrección cambió el costo: ahora usa los pasajeros asignados a cada molinete (los que caminan hacia
él y los que esperan) multiplicados por el tiempo de validación de referencia, que se conoce antes del
sorteo. En la misma revisión se corrigieron otras fallas menores:
\begin{itemize}
  \item cuatro arribos fantasma, de una lanzadera que no llega a Constitución y que se colaba al leer el
        horario;
  \item la espera de quien pasaba directo, que incluía unos 32~s de caminata;
  \item la cola máxima, que se medía también durante el drenaje;
  \item el fin técnico de la corrida, que se fijó en \num{17900}~s para emitir la fila antes del tope de
        5~horas.
\end{itemize}

Con valores de prueba, el nuevo piloto de \EO{} drenó los \num{18545} pasajeros. Los 20 molinetes
quedaron en uso, con utilizaciones de entre 17 y 52~\%. El caudal simulado por ventana quedó en su
mayoría a menos del 6~\% del observado. Las mayores diferencias fueron en las ventanas de 07:15
($-16$~\%) y 08:00 ($+15$~\%). La corrida tardó 274~s con la animación visible y 28~s sin ella.
El par \EO{}-\EI{} con la misma semilla generó exactamente la misma demanda, lo que confirmó en el motor
los números aleatorios comunes.

\subsection{Un peatón encerrado}

Con el escenario intermedio ya cargado, dos intentos de producción se cortaron: uno a las 22 corridas y
otro a las 7. Un peatón
empujado por la multitud quedaba dentro del contorno cerrado de la cabina o de los cierres de los bancos
de molinetes y no podía salir, y la corrida terminaba sin drenar. Esos contornos se reemplazaron por
paredes abiertas. Las filas de esos dos intentos se descartaron: se conservan aparte y no se cargaron en
la planilla.

\subsection{Tiempo de cómputo}

La tabla~\ref{tab:computo} resume la evolución del tiempo por corrida. La caída de 43~minutos a segundos
no se debe a optimizar el código: el sistema dejó de atascarse. Un vestíbulo que drena tiene pocos
peatones activos a la vez, y el costo de la simulación peatonal crece con esa cantidad.

\begin{table}[H]
\centering
\caption{Tiempo real por corrida de la franja completa}
\label{tab:computo}
\small
\begin{tabular}{@{}llr@{}}
\toprule
Piloto & Condición & Tiempo \\
\midrule
24/09, \EO{}, valores de prueba & Atascado hasta el tope de 5~h & $\approx 43$~min \\
05/10, \EO{}, valores de prueba & Drena; con animación & 274~s \\
05/10, \EI{}, valores de prueba & Drena; sin animación & 28~s \\
05/10, producción, escenario intermedio & Drena; sin animación & $\approx 7$~s \\
\bottomrule
\end{tabular}
\end{table}

Con 7~s por corrida, los 30 pares de producción llevan unos 7 minutos. El cómputo dejó de ser una
restricción del diseño experimental.

\subsection{Qué se aprendió}

Las tres fallas tienen algo en común: ninguna producía un error, solo resultados plausibles pero
equivocados o corridas que no terminaban. La elección de cola sin servicio habría dado esperas exageradas
y un \EI{} artificialmente bueno. Las encontraron dos cosas: instrumentar lo que las métricas no miraban
(el hall) y comparar contra el ejemplo oficial de la herramienta. Por eso el verificador incorporó un
control para cada una:
\begin{itemize}
  \item la oleada se reparte entre colas aun sin servicio sorteado;
  \item la espera de quien pasa directo es cero;
  \item las colas tienen la orientación correcta.
\end{itemize}
